\chapter{TÓM LƯỢC}

\indent\indent Trong bối cảnh bùng nổ của các mạng xã hội phi tập trung (như Lens Protocol), việc phát hiện các cuộc tấn công Sybil trở thành thách thức sống còn để bảo vệ tính toàn vẹn và sự tin cậy của hệ sinh thái. Các phương pháp truyền thống thường gặp khó khăn trong việc khai thác các mối quan hệ ràng buộc đa tầng phức tạp và đối mặt với sự thiếu hụt nghiêm trọng dữ liệu gán nhãn chuẩn. Nghiên cứu này đề xuất \textbf{SybilSignal}, một mô hình đồ thị trí tuệ nhận thức ràng buộc hỗ trợ phát hiện Sybil có khả năng giải thích. Hệ thống mô hình hóa dữ liệu xã hội đa chiều thành cấu trúc đồ thị quan hệ 4 tầng (bao gồm Profile, Interact, similarity và Co-owner) kết hợp với cơ chế chuẩn hóa trọng số logarit nhằm xử lý nhiễu trong các mạng lưới có mật độ kết nối cao. Nghiên cứu áp dụng kiến trúc Decoupled GNN-ML, sử dụng mạng chú ý đồ thị cải tiến (GATv2) để trích xuất đặc trưng kết hợp với các thuật toán học máy quần thể để phân lớp. Kết quả thực nghiệm cho thấy cấu hình GATv2 + Random Forest đạt chỉ số Macro F1-Score vượt trội là 0.9303 trên tập dữ liệu quy mô lớn theo phạm vi Quý. Đặc biệt, mô hình đạt tỷ lệ nhận diện nhầm người dùng thật là 0\%, khẳng định độ an toàn và tin cậy cao khi triển khai thực tế. Với thiết kế tách rời giúp tối ưu hóa đột phá chi phí tính toán và tăng cường tính minh bạch, nghiên cứu cung cấp một giải pháp hiệu quả và bền vững cho công tác quản trị an ninh trong môi trường Web3.\vspace{0.3cm}

\textbf{Từ khóa:} Phát hiện Sybil, Đồ thị xã hội Web3, Mạng nơ-ron đồ thị, GATv2, Kiến trúc Decoupled, Random Forest, SybilSignal.